\documentclass[12pt,a4paper]{article}

\usepackage[utf8]{inputenc}
\usepackage[T1]{fontenc}
\usepackage[french]{babel}
\usepackage{geometry}
\usepackage{hyperref}
\usepackage{xcolor}
\usepackage{listings}
\usepackage{enumitem}
\usepackage{amsmath}

\geometry{margin=2.5cm}

\lstdefinestyle{cppstyle}{
    language=C++,
    basicstyle=\ttfamily\small,
    keywordstyle=\color{blue},
    commentstyle=\color{green!40!black},
    stringstyle=\color{red!60!black},
    numbers=left,
    numberstyle=\tiny,
    stepnumber=1,
    numbersep=8pt,
    showstringspaces=false,
    breaklines=true,
    frame=single
}

\title{Rapport TP1 – Troisi\`eme Partie (code)}
\author{Nom : \`a compl\'eter \\ Fili\`ere : \`a compl\'eter}
\date{\today}

\begin{document}

\maketitle
\tableofcontents
\newpage

\section{Exercice 16 – Calculer l'intersection de plusieurs rectangles}
\paragraph{}
L'exercice 16 repose sur une unique fonction bool\'eenne \texttt{Rectangle::intersection} associ\'ee \`a une boucle d'intersection successive. Ce rapport met l'accent sur le code qui rend la d\'emonstration concr\`ete : la d\'etection d'un chevauchement et la r\'eduction d'un tableau de rectangles \`a un seul objet commun.

\subsection{Intersection bool\'eenne}
La fonction suivante compare les axes x et y et ne retourne une intersection que si les bornes se chevauchent.
\begin{lstlisting}[style=cppstyle]
bool Rectangle::intersection(const Rectangle& r, Rectangle& resultat) const {
    const int ix1 = std::max(x1_, r.x1_);
    const int iy1 = std::max(y1_, r.y1_);
    const int ix2 = std::min(x2_, r.x2_);
    const int iy2 = std::min(y2_, r.y2_);

    if (ix1 > ix2 || iy1 > iy2) {
        return false;
    }

    resultat = Rectangle(ix1, iy1, ix2, iy2);
    return true;
}
\end{lstlisting}

\paragraph{}
Cette version bool\'eenne supprime l'ambigu\"ite du tuple $(0,0,0,0)$ et sert de fondement \`a la boucle suivante.

\subsection{Intersection d'une collection}
On applique la fonction pr\'ec\'edente sur chaque rectangle d'un tableau pour conserver l'intersection courante.
\begin{lstlisting}[style=cppstyle]
void afficherIntersectionsMultiples(const std::vector<Rectangle>& rectangles) {
    if (rectangles.empty()) {
        std::cout << "Aucun rectangle saisi.\n";
        return;
    }

    Rectangle courant = rectangles.front();
    for (size_t i = 1; i < rectangles.size(); ++i) {
        Rectangle resultat;
        if (!courant.intersection(rectangles[i], resultat)) {
            std::cout << "Intersection globale vide apres le rectangle " << i + 1 << ".\n";
            return;
        }
        courant = resultat;
    }

    std::cout << "Intersection globale: ";
    courant.afficher(std::cout);
    std::cout << '\n';
}
\end{lstlisting}

\paragraph{}
La boucle s'arrête dès qu'une intersection est vide, ce qui respecte la sémantique attendue : l'intersection globale n'existe plus.

\subsection{Structure complète de \texttt{Rectangle}}
Le cadrage précis de toutes les méthodes pertinentes rappelle que l'exercice tire parti des fonctions existantes.
\begin{lstlisting}[style=cppstyle]
class Rectangle {
public:
    Rectangle(int x1 = 0, int y1 = 0, int x2 = 0, int y2 = 0);
    Rectangle(const Point& p1, const Point& p2);

    void afficher(std::ostream& os) const;
    void saisir(std::istream& is);

    bool carre() const;
    int surface() const;
    void rotation90();
    int position(const Point& p) const;
    void translation(const Point& p);
    int comparer(const Rectangle& r) const;
    int position(const Rectangle& r) const;
    Rectangle reunion(const Rectangle& r) const;
    Rectangle intersection(const Rectangle& r) const;
    bool intersection(const Rectangle& r, Rectangle& resultat) const;

private:
    void ordonner();
    int x1_;
    int y1_;
    int x2_;
    int y2_;
};
\end{lstlisting}

\section{Exercice 17 – Rectangle définie par deux points}
\paragraph{}
L'exercice 17 demande de recomposer la classe \texttt{Rectangle} en utilisant des objets \texttt{Point}. La solution garde la même logique qu'auparavant, mais change la représentation interne pour illustrer la composition.

\subsection{Déclaration de \texttt{RectanglePoints}}
\begin{lstlisting}[style=cppstyle]
class RectanglePoints {
public:
    RectanglePoints(const Point& hautGauche = Point(), const Point& basDroite = Point());

    void afficher(std::ostream& os) const;
    int surface() const;
    bool carre() const;
    int position(const Point& p) const;
    RectanglePoints reunion(const RectanglePoints& r) const;
    bool intersection(const RectanglePoints& r, RectanglePoints& resultat) const;

private:
    void ordonner();

    Point hautGauche_;
    Point basDroite_;
};
\end{lstlisting}

\paragraph{}
Cette interface conserve l'ensemble des opérations du rectangle à quatre entiers, mais délègue la gestion des coordonnées aux objets \texttt{Point}.

\subsection{Ordonnancement et intersection}
\begin{lstlisting}[style=cppstyle]
void RectanglePoints::ordonner() {
    const int xMin = std::min(hautGauche_.abscisse(), basDroite_.abscisse());
    const int yMin = std::min(hautGauche_.ordonnee(), basDroite_.ordonnee());
    const int xMax = std::max(hautGauche_.abscisse(), basDroite_.abscisse());
    const int yMax = std::max(hautGauche_.ordonnee(), basDroite_.ordonnee());

    hautGauche_.setAbscisse(xMin);
    hautGauche_.setOrdonnee(yMin);
    basDroite_.setAbscisse(xMax);
    basDroite_.setOrdonnee(yMax);
}

bool RectanglePoints::intersection(const RectanglePoints& r, RectanglePoints& resultat) const {
    const int ix1 = std::max(hautGauche_.abscisse(), r.hautGauche_.abscisse());
    const int iy1 = std::max(hautGauche_.ordonnee(), r.hautGauche_.ordonnee());
    const int ix2 = std::min(basDroite_.abscisse(), r.basDroite_.abscisse());
    const int iy2 = std::min(basDroite_.ordonnee(), r.basDroite_.ordonnee());

    if (ix1 > ix2 || iy1 > iy2) {
        return false;
    }

    resultat = RectanglePoints(Point(ix1, iy1), Point(ix2, iy2));
    return true;
}
\end{lstlisting}

\paragraph{}
Les opérations d'ordonnancement et d'intersection s'appuient sur les accesseurs/mutateurs de `Point`, montrant que le comportement global reste identique malgré la nouvelle représentation.

\end{document}
